% ═══════════════════════════════════════════════════════════════════════
\section{Le théorème du point fixe de Banach}
% ═══════════════════════════════════════════════════════════════════════

\begin{frame}{Théorème 2.3 --- Théorème du point fixe de Banach}
  \small Complétude et contraction : tous les ingrédients sont en place.
  On énonce et démontre d'abord le théorème dans $(\R^2,\dE)$, pour
  pouvoir dessiner chaque étape.
  \begin{Enonce}{Théorème 2.3 --- Théorème du point fixe de Banach}
    Soit $f:\R^2\to\R^2$ une contraction de rapport $k$ (pour $\dE$).
    Alors :
    \begin{enumerate}
      \item[(i)] $f$ possède un unique point fixe
        $\overline{x}\in\R^2$ ;
      \item[(ii)] pour tout $x_0\in\R^2$, la suite $x_n = f^n(x_0)$
        converge vers $\overline{x}$, quel que soit $x_0$ ;
      \item[(iii)] on a la majoration
        $\displaystyle \dE(x_n,\overline{x})\leq\frac{k^n}{1-k}\,\dE(x_0,x_1)$.
    \end{enumerate}
  \end{Enonce}
  \scriptsize ($f^n = f\circ f\circ\cdots\circ f$, $n$ fois : une
  composée, pas un produit.)
\end{frame}

\begin{frame}{Le plan de la démonstration}
  \FigBox[0.5\textheight]{%
    \begin{tikzpicture}[
        etape/.style={draw=sectioncolor, fill=thmboxbg, rounded corners=2pt,
          thick, align=center, font=\scriptsize, minimum height=1.25cm,
          text width=2.3cm}]
      \node[etape] (a) at (0,0) {Partir d'un $x_0$\\quelconque ;\\$x_n = f^n(x_0)$};
      \node[etape] (b) at (3.0,0) {$f^n$ contraction\\de rapport $k^n$\\(lemme 2.4)};
      \node[etape] (c) at (6.0,0) {$(x_n)$ est\\de Cauchy\\(lemme 2.5)};
      \node[etape, fill=secboxbg, draw=coveraccent!85!black] (d) at (9.0,0)
        {$\R^2$ complet :\\$x_n\to\overline{x}$\\(remarque 2.8)};
      \node[etape] (e) at (1.5,-2.0) {$f(\overline{x})=\overline{x}$\\(lemme 2.6)};
      \node[etape] (f) at (4.5,-2.0) {Majoration\\de l'erreur\\(prop.~2.7)};
      \node[etape] (g) at (7.5,-2.0) {Unicité\\(lemme 2.8)};
      \draw[fleche] (a) -- (b);
      \draw[fleche] (b) -- (c);
      \draw[fleche] (c) -- (d);
      \draw[fleche] (d.south) to[out=-90, in=0] ($(d.south)+(-0.6,-0.6)$)
        -| (e.north);
      \draw[fleche] (e) -- (f);
      \draw[fleche] (f) -- (g);
    \end{tikzpicture}}
  \small Seul ingrédient propre à $\R^2$ : sa complétude.
  Tout le reste n'utilise que les axiomes d'une métrique.
\end{frame}

\begin{frame}{Lemme 2.4 --- Une contraction itérée reste une contraction}
  \begin{columns}[c]
    \begin{column}{0.56\textwidth}
      \begin{Enonce}{Lemme 2.4}
        \small Soit $f:\R^2\to\R^2$ une contraction de rapport $k$.
        Pour tout $n\in\Ns$, $f^n$ est une contraction de rapport $k^n$ :
        \[
          \dE\big(f^n(x),f^n(y)\big)\leq k^n\,\dE(x,y).
        \]
      \end{Enonce}
      \begin{AuTableau}
        \footnotesize Récurrence sur $n$ : appliquer la contraction aux
        points $f^n(x)$, $f^n(y)$, puis l'hypothèse de récurrence.
      \end{AuTableau}
    \end{column}
    \begin{column}{0.42\textwidth}
      \FigBox[0.6\textheight]{%
        \begin{tikzpicture}[scale=0.95]
          \node[pt, label={[etiq]left:$x$}] (x0) at (0,3.6) {};
          \node[pt, label={[etiq]right:$y$}] (y0) at (4.6,3.6) {};
          \draw[sectioncolor, thick] (x0) -- (y0);
          \node[etiq, sectioncolor, above] at (2.3,3.6) {$\dE(x,y)$};
          \node[pt, label={[etiq]left:$f(x)$}] (x1) at (1.15,2.1) {};
          \node[pt, label={[etiq]right:$f(y)$}] (y1) at (3.45,2.1) {};
          \draw[sectioncolor, thick] (x1) -- (y1);
          \node[etiq, sectioncolor, above] at (2.3,2.1) {$k\,\dE(x,y)$};
          \node[pt, label={[etiq]below:$f^2(x)$}] (x2) at (1.6,0.6) {};
          \node[pt, label={[etiq]below:$f^2(y)$}] (y2) at (3.0,0.6) {};
          \draw[sectioncolor, thick] (x2) -- (y2);
          \node[etiq, sectioncolor, above] at (2.3,1.05) {$k^2\dE(x,y)$};
          \draw[fleche, sectioncolor!60] (x0) -- (x1);
          \draw[fleche, sectioncolor!60] (y0) -- (y1);
          \draw[fleche, sectioncolor!60] (x1) -- (x2);
          \draw[fleche, sectioncolor!60] (y1) -- (y2);
        \end{tikzpicture}}
      \begin{center}\scriptsize Figure 2.14\end{center}
    \end{column}
  \end{columns}
\end{frame}

\begin{frame}{Lemme 2.5 --- La suite des itérées est de Cauchy}
  \begin{Enonce}{Lemme 2.5}
    \small Soit $f$ une contraction de rapport $k$, $x_0\in\R^2$ et
    $x_n=f^n(x_0)$. Pour tous $n,p\in\N$,
    \[
      \dE(x_n,x_{n+p})\leq\frac{k^n(1-k^p)}{1-k}\,\dE(x_0,x_1),
      \tag{2.2}
    \]
    et par conséquent $(x_n)$ est de Cauchy dans $(\R^2,\dE)$.
  \end{Enonce}
  \begin{AuTableau}
    \footnotesize Lemme 2.4 : $\dE(x_i,x_{i+1})\leq k^i\dE(x_0,x_1)$.
    Inégalité triangulaire $p$ fois le long de la ligne brisée
    (figure 2.15), puis somme géométrique. Comme $1-k^p\leq1$ et
    $k^n\to0$, la borne tend vers $0$ uniformément en $p$.
  \end{AuTableau}
\end{frame}

\begin{frame}{Figure 2.15 --- La ligne brisée}
  \FigBox[0.36\textheight]{%
    \begin{tikzpicture}[scale=1.2]
      \node[pt, label={[etiq]below:$x_n$}] (xn) at (0,0) {};
      \node[pt, label={[etiq]above:$x_{n+1}$}] (xn1) at (1.3,0.75) {};
      \node[pt, label={[etiq]above:$x_{n+2}$}] (xn2) at (2.35,1.05) {};
      \node[etiq] (dots) at (3.35,0.7) {$\cdots$};
      \node[pt, label={[etiq]below:$x_{n+p-1}$}] (xnp1) at (4.35,0.35) {};
      \node[pt, label={[etiq]above:$x_{n+p}$}] (xnp) at (5.6,1.1) {};
      \draw[sectioncolor, thick] (xn) -- (xn1) -- (xn2);
      \draw[sectioncolor, thick, dashed] (xn2) -- (xnp1);
      \draw[sectioncolor, thick] (xnp1) -- (xnp);
      \draw[firebrick3, thick, dashed] (xn) --
        node[midway, below, etiq, firebrick3]{$\dE(x_n,x_{n+p})$} (xnp);
    \end{tikzpicture}}
  \small
  \[
    \dE(x_n,x_{n+p}) \leq \sum_{i=0}^{p-1}\dE(x_{n+i},x_{n+i+1})
    \leq \dE(x_0,x_1)\sum_{i=0}^{p-1}k^{n+i}.
  \]
  La diagonale (en rouge) d'un polygone, même croisé, n'excède jamais
  son périmètre.
\end{frame}

\begin{frame}{Lemme 2.6 --- La limite est un point fixe}
  \small $(\R^2,\dE)$ est complet et $(x_n)$ est de Cauchy : il existe
  $\overline{x}\in\R^2$ avec $x_n\to\overline{x}$.
  \begin{Enonce}{Lemme 2.6}
    Avec les notations précédentes, $f(\overline{x})=\overline{x}$.
  \end{Enonce}
  \begin{AuTableau}
    \small Une contraction est continue ; elle commute donc avec la
    limite : $f(\lim x_n) = \lim f(x_n) = \lim x_{n+1} = \overline{x}$.
  \end{AuTableau}
\end{frame}

\begin{frame}{Proposition 2.7 --- Majoration de l'erreur}
  \begin{Enonce}{Proposition 2.7 --- Majoration de l'erreur}
    Avec les notations précédentes,
    \[
      \dE(x_n,\overline{x})\leq\frac{k^n}{1-k}\,\dE(x_0,x_1)
      \qquad\text{pour tout } n\in\N.
    \]
  \end{Enonce}
  \begin{AuTableau}
    \small Faire $p\to\infty$ dans (2.2) ; continuité de
    $\dE(x_n,\cdot)$ et $k^p\to0$.
  \end{AuTableau}
  \vspace{4pt}
  \small Usage pratique : on sait à l'avance combien
  d'itérations suffisent pour une précision voulue, à partir du seul
  premier pas $\dE(x_0,x_1)$.
\end{frame}

\begin{frame}{Lemme 2.8 --- Unicité du point fixe}
  \begin{Enonce}{Lemme 2.8 --- Unicité du point fixe}
    Une contraction $f:\R^2\to\R^2$ possède au plus un point fixe.
  \end{Enonce}
  \begin{AuTableau}
    \small Par l'absurde : deux points fixes distincts donnent
    $\dE(\overline{x}_1,\overline{x}_2)\leq k\,\dE(\overline{x}_1,\overline{x}_2)$ ;
    on divise (séparation !) et on obtient $1\leq k$.
  \end{AuTableau}
  \vspace{6pt}
  \small Lemmes 2.4, 2.5, 2.6, 2.8 et proposition 2.7 réunis :
  toutes les affirmations du théorème 2.3 sont établies.
\end{frame}

\begin{frame}{Figure 2.16 --- Interprétation géométrique}
  \begin{columns}[c]
    \begin{column}{0.45\textwidth}
      \small À chaque application de $f$, la distance au point fixe est
      multipliée par un facteur $k<1$.
      \vspace{6pt}

      La majoration montre que $x_n\to\overline{x}$ quel que
      soit $x_0$ --- et toujours vers le même point fixe,
      peu importe où on commence à itérer.
    \end{column}
    \begin{column}{0.53\textwidth}
      \FigBox[0.6\textheight]{%
        \begin{tikzpicture}[scale=1.7]
          \draw[fig] (-0.4,-0.6) rectangle (3.6,2.4);
          \node[etiq, sectioncolor] at (0.1,2.0) {$\R^2$};
          \node[pt, label={[etiq]below left:$x_0$}] (x0) at (0.5,0.35) {};
          \node[pt, label={[etiq]below:$x_1$}] (x1) at (1.8,0.95) {};
          \node[pt, label={[etiq]right:$x_2$}] (x2) at (2.3,1.45) {};
          \node[ptim, label={[etiq]above:$\overline{x}$}] (xbar) at (2.6,1.75) {};
          \draw[fleche, sectioncolor] (x0) to[bend left=15] node[midway, below right, etiq]{$f$} (x1);
          \draw[fleche, sectioncolor] (x1) to[bend left=15] node[midway, above left, etiq]{$f$} (x2);
          \draw[sectioncolor!50, thick, dotted] (x2) to[bend left=15] (xbar);
        \end{tikzpicture}}
    \end{column}
  \end{columns}
\end{frame}

\begin{frame}{Remarque 2.11 --- Le cas d'un espace métrique quelconque}
  \begin{Rem}{Remarque 2.11}
    \small Aucun argument des lemmes 2.4, 2.5, 2.6, 2.8 ni de la
    proposition 2.7 n'utilise une propriété particulière du plan :
    seulement les quatre axiomes d'une métrique et la
    complétude.
  \end{Rem}
  \begin{Enonce}{Théorème de Banach --- forme générale}
    \small Si $(E,d)$ est un espace métrique complet et
    $f:E\to E$ une contraction, alors $f$ possède un unique point fixe
    $\overline{x}$, vers lequel converge $(f^n(x_0))$ pour tout
    $x_0\in E$, avec la même majoration de l'erreur.
  \end{Enonce}
  \small Deux usages à venir : un intervalle fermé de $\R$ (méthode de
  Newton, tout de suite) ; l'espace $\big(\Comp{\R^2},\dH\big)$
  (fractales, chapitre 3).
\end{frame}
